% Ch. 6 : segmentation en zones
\documentclass[border=6pt]{standalone}
\usepackage{coursfig}
\begin{document}
\begin{tikzpicture}[x=1mm,y=1mm]
  \node[box=Ocre, text width=40mm, minimum height=22mm] (lb) at (0,0)
       {\textbf{listify-lb}\sub{seul joignable de}\sub{l'extérieur (80, 443)}};
  \node[box=Enc, text width=30mm] (a1) at (80,12) {\textbf{listify-app1}};
  \node[box=Enc, text width=30mm] (a2) at (80,-12) {\textbf{listify-app2}};
  \node[db=Plum, minimum width=30mm] (db) at (160,0) {\textbf{listify-db}};

  \foreach \a in {a1,a2} {
    \draw[flow] (lb.east) -- (\a.west);
    \draw[flow=Plum] (\a.east) -- (db.west);
  }
  \node[elabelc, fill=none] at (42,16) {8000};
  \node[elabel, fill=none, text=Muted] at (42,-16) {uniquement};
  \node[elabelc, fill=none] at (122,16) {5432};
  \node[elabel, fill=none, text=Muted] at (122,-16) {uniquement};

  \begin{scope}[on background layer]
    \node[dframe=Ocre, fit=(lb), inner sep=5mm, yshift=3mm, inner ysep=8mm] (z1) {};
    \node[dframe=Enc,  fit=(a1)(a2), inner sep=5mm, yshift=3mm, inner ysep=8mm] (z2) {};
    \node[dframe=Plum, fit=(db), inner sep=5mm, yshift=3mm, inner ysep=8mm] (z3) {};
  \end{scope}
  \node[ftitle, text=Ocre, anchor=north] at (z1.north) {ZONE EXPOSÉE};
  \node[ftitle, text=Enc,  anchor=north] at (z2.north) {ZONE APPLICATIVE};
  \node[ftitle, text=Plum, anchor=north] at (z3.north) {ZONE DONNÉES};

  \node[box=Brick, fill=BrickBg, text width=30mm] (x) at (0,-46) {tout autre flux};
  \draw[dflow=Brick, -] (x.east) -| (z3.south);
  \node[circle, fill=Brick, text=white, font=\sffamily\bfseries, inner sep=0pt, minimum size=5mm] at (z3.south) {$\times$};
  \node[note, text=Brick, above=1mm] at (80,-46) {refusé par les pare-feu de chaque machine};
\end{tikzpicture}
\end{document}
